\chapter{ Representation Theory}\label{MSrep}

In this  chapter we introduce the integral equations that govern the 
time harmonic elastic wavefield  
in an arbitrary $3-$D isotropic elastic halfspace with superimposed
$2-$D space filling heterogeneity. The formulation is abstract, because we 
do not provide {\em constructive representations\/} for the various  
operators, introduced here. 
Computational representations for these operators will be developed 
incrementally in the later chapters of this thesis. The present chapter 
serves as a prelude to those developments.

Our presentation differs from previous work, in that it is based 
directly on the first order displacement-traction equations rather than the 
more common coupled second order displacement field equations. 
We derive a {\em displacement-traction representation theorem\/} in 
\S\ref{section2Drep} that plays a role analogous to the {\em Rayleigh-Betti representation 
theorem\/} of the corresponding coupled second order theory. We then 
employ the new representation theorem to derive the time harmonic {\em 
elastic wave Lippmann-Schwinger type integral equation \/} in \S\ref{secLS}.




 
\section{The multiple scattering problem}

\subsection{The multiple scattering problem geometry}

\begin{figure}
\HideDisplacementBoxes
%\ShowDisplacementBoxes
\centerline{\BoxedEPSF{MSRep1.EPSF scaled 1000}}
\caption{The Representation Theorem Geometry}
\protect\label{MSRep1}
\end{figure}

We partition the plane ${\bmi x}=(z,x)\in{ R}^{2}$ into
the disjoint union of the {\em scattering zone} ${\mit \Omega}$ of compact support and its
complement ${\mit \Omega}^{e}$ which we will
refer to as the {\em external zone} (Fig.\ref{MSRep1}). The oriented 
boundary of ${\mit \Omega}$ is denoted $\partial {\mit \Omega}^{+}$, 
while $\partial {\mit \Omega}^{\infty}$ is the {\em virtual semi-circular 
boundary\/} at infinity, with $\partial {\mit \Omega}^{f}$ the {\em free surface\/} component of 
the external boundary.  We 
cut ${\mit \Omega}^{e}$ along $\vartheta$, to ensure that ${\mit \Omega}^{e}$ is 
{\em simply connected}. The contribution from the cut $\vartheta$ 
cancels because the wavefield is continuous across the cut. The 
cuts ensure  that the boundary contour 
$\partial{\mit \Omega}^{+}$ of the scattering zone ${\mit \Omega}$ has 
the proper sense. 

Let $\chi_{{\mit \Omega}}$ denote the {\em characteristic function } on ${{\mit \Omega}}$,
defined by 
\begin{equation}
\chi_{{\mit \Omega}}=
\left\{
\begin{array}{ccl} 1&:&{\bmi x}\in {\mit \Omega}\\ 0&:&{\bmi x}\not\in {\mit \Omega}
\end{array}  
\right. ,
\end{equation} 
so that $\chi_{{\mit \Omega}}$ takes unit value in the scattering zone
but identically vanishes in the external zone. The isotropic material properties ${\bmi
m}(z,x)=(\rho,\alpha,\beta)(z,x)$, which denote the density, 
compressional and shear-wavespeeds respectively, are
defined as the sum of two parts: a {\em reference} component
${\rf {\bmi m}}=({\rf \rho},{\rf \alpha},{\rf \beta})$, 
defined throughout the halfspace; and a {\em scattering} component 
${\mit \Delta} {\bmi m}(z,x)=({\mit \Delta}{\rho},{\mit \Delta}{\alpha},{\mit \Delta}{\beta})$ 
confined to the  scattering zone. Hence,  
\begin{equation} \label{mdecomp}
{\bmi m}(z,x)=
{\rf {\bmi m}}+\chi_{{\mit \Omega}} {\mit \Delta}{\bmi m}(z,x).
\end{equation}

\subsection{The Governing equations}


Assuming that the depth below the free surface is measured by the 
$z-$coordinate, the forced elastic displacement ${\bmi  u} = {({u}_{z},{u}_{x},{u}_{y})}^{T}$
and traction vector ${\bmi  t} = {({\tau}_{zz},{\tau}_{xz},{\tau}_{yz})}^{T }$ on a 
fixed depth plane (i.e. $z=const.$) with unit
normal ${\bmi  e}_{z}$,  are grouped into the $6$-vector
\begin{eqnarray} \label{bdef}
\lefteqn{{\bmi  b}  (z,x,y)= {  \left[ \begin{array}{cccccc}{ u}_{ z}
(z,x,y),&{u}_{x}(z,x,y),&{u}_{y}(z,x,y),&{\tau}_{zz}(z,x,y),&{\tau}_{xz}(z,x,y),&
{\tau}_{yz}(z,x,y)\end{array} \right]}^{  T }.}\nonumber\\
&&
\end{eqnarray}
The governing equation for time harmonic elastic wave motion in an {\em isotropic  medium\/} 
is then a 
set of $6$ coupled partial differential equations (first order with respect 
to the depth coordinate) which may be written in the operator notation
\begin{equation} \label{govEqn00} 
{\CB L}({\bmi x};{\partial }_{z},{\partial }_{x},{\partial }_{y})\cdot {\bmi  b }\equiv\left[\partial_{z}-{\CB A}({\bmi x};{\partial }_{x},{\partial }_{y})\right]\cdot {\bmi  b }=
{\bmi  f}
\end{equation}
where ${\bmi  b}$ is given in (\ref{bdef})  and 
${\CB A}$ is defined as
\begin{eqnarray} \label{Adef0}
\lefteqn{{\CB A}({\bmi x};{\partial }_{x},{\partial }_{y})=}\nonumber\\
&&\left[{\begin{array}{cccccc}
0&- \gamma {\partial }_{  x}&  - \gamma {\partial }_{  y}&  a&0&0\\ 
-{ \partial }_{x}&0&0&0&b&0\\ 
-{\partial }_{y}&0&0&0&0&b\\ 
- \rho {\omega }^{  2}&  0&0&0&-{ \partial }_{x}&-{ \partial }_{y}\\ 
0&- \rho {\omega }^{  2}  -{\partial }_{x} \zeta {\partial }_{  x}  -{ \partial }_{y} \mu 
{\partial }_{  y}&  -{ \partial }_{x} \chi {\partial }_{  y} 
 -{ \partial }_{y} \mu {\partial }_{  x}&  -{ \partial }_{x} \gamma &  0&0\\
 0&-{ \partial }_{y} \chi {\partial }_{  x}  -{ \partial }_{x} \mu {\partial }_{  y}&
   - \rho {\omega }^{  2}  -{ \partial }_{y} \zeta {\partial }_{  y} 
 -{ \partial }_{x} \mu {\partial }_{  x}&  -{ \partial }_{y} \gamma &  0&0\end{array}}
 \right]
\end{eqnarray} 
and the material parameters in (\ref{Adef0}) are defined in Table \ref{ATable}.

%%%%%%%%%%%%%%%%%%%%%%%%%%%%%%%%%%%%%%%%%%%%%%%%%%%%%%%%%%%%%%%%%%%%%%%%%%%%%%%%%%%%%%
%%%%%%%%%%%%%%%%%%%%%%%%%%%%%%%%%%%%%%%%%%%%%%%%%%%%%%%%%%%%%%%%%%%%%%%%%%%%%%%%%%%%%%
%%%%%%%%%%%%%%%%%%%%%%%%%%%%%%%%%%%%%%%%%%%%%%%%%%%%%%%%%%%%%%%%%%%%%%%%%%%%%%%%%%%%%%
\begin{table}
\begin{tabular*}{85mm}{ccc}\hline
parameter & $(\rho, \alpha, \beta)$-definition & $(\rho, \lambda, \mu)$-definition\\ \hline
$P$&$\alpha$&$\sqrt{{{\lambda + 2 \mu} \over {\rho}}}$\\
$S$&$\beta$&$\sqrt{{{\mu} \over {\rho}}}$\\
$H$&$\beta$&$\sqrt{{{\mu} \over {\rho}}}$\\
$\gamma$&$\left(1- 2{{\beta^2}\over{\alpha^2}}\right)$&${{\lambda}\over{\lambda+2 \mu}}$\\ 
$\zeta$&$4 \rho \beta^2 \left(1- {{\beta^2}\over{\alpha^2}} \right)$&${{4 
\mu (\lambda+ \mu)}\over{\lambda+2 \mu}}$\\ 
$\chi$&$2 \rho \beta^2\left(1- 
2{{\beta^2}\over{\alpha^2}}\right)$&${{2 \lambda \mu}\over{\lambda+2 
\mu}}$\\ $a$&${{1}\over{\rho \alpha^2}}$&${{1}\over{\lambda+2 \mu}}$\\ 
$b$&${{1}\over{\rho \beta^2}}$&${{1}\over{\mu}}$\\ 
\hline
\end{tabular*}
\caption{\label{ATable} { The material parameter definitions}}
\end{table} 

This thesis is concerned with finding solutions of (\ref{govEqn00}) when 
the material parameters in Table \ref{ATable} are functions of $(z,x)$ but are 
independent of  the $y-$coordinate. This problem is sometimes called the 
$2 \slantfrac{1}{2}-$D problem because the wavefield may spread in $3-$D, 
but the heterogeneity is confined to $2-$D.
We exploit the translational invariance of the medium along the 
$y-$coordinate by assuming that the displacement/traction field ${\bmi 
b}(z,x,y)$ may be {\em synthesized\/} from its  $k_y$ plane wave components as
\begin{equation} \label{kysyn}
{\bmi b}(z,x,y)={\frac{1}{2\pi }}\int_{R} \, {\rm d}{k}_{y}\,{\bmi 
b}(z,x,{k}_{y})\, {\rm e}^{i{k}_{y}y}
\end{equation}
where ${\bmi b}(z,x,{k}_{y})$ is now governed by
\begin{equation} \label{govEqn0} 
{\CB L}(z,x;{\partial }_{z},{\partial }_{x},i {k}_{y})\cdot {\bmi  b 
}(z,x,k_y; \omega)\equiv
\left[\partial_{z}-{\CB A}(z,x;{\partial }_{x},i {k}_{y})\right]\cdot {\bmi  b }(z,x,k_y; 
\omega)=
\tilde{\bmi  f}
\end{equation}
with
\begin{eqnarray} \label{Adef}
\lefteqn{{\CB A}(z,x;{\partial }_{x},i {k}_{y})
=}\nonumber\\
&&\left[{\begin{array}{cccccc}
0&- \gamma {\partial }_{  x}&  -i \gamma {  k}_{  y}&  a&0&0\\ -{ \partial 
}_{x}&0&0&0&b&0\\ -i{k}_{y}&0&0&0&0&b\\ - \rho {\omega }^{  
2}&  0&0&0&-{ \partial }_{x}&-i{k}_{y}\\ 0&- \rho {\omega 
}^{  2}  -{ \partial }_{x} \zeta {\partial }_{  
x}  + \mu {  k}_{  y}^{  2}&  
-i{k}_{y}{ \partial }_{x} \chi   -i{k}_{y} \mu {\partial 
}_{  x}&  -{ \partial }_{x} \gamma &  0&0\\ 
0&-i{k}_{y} \chi {\partial }_{  x}  -i{k}_{y}{ \partial 
}_{x} \mu &  - \rho {\omega }^{  2}  + \zeta 
{  k}_{  y}^{  2}  -{ \partial }_{x} \mu 
{\partial }_{  x}&  -i{k}_{y} \gamma &  0&0\end{array}}\right]
\end{eqnarray}
and
\begin{equation} \label{kysynf}
\tilde{\bmi f}=\int_{R} \, {\rm d}{k}_{y}\,{\bmi f}(z,x,y)
\, {\rm e}^{-i{k}_{y}y}.
\end{equation}
In this work, we assume the the motion is driven by a point source in the 
$y=0$ plane, i.e. ${\bmi f}(z,x,y)= {\bmi F}(z,x) \delta(y)$, so that 
from (\ref{kysynf}) we have 
\begin{equation} \label{kysynf1}
\tilde{\bmi f}=\int_{R} \, {\rm d}{k}_{y}\,{\bmi F}(z,x) \delta(y)
\, {\rm e}^{-i{k}_{y}y} = {\bmi F}(z,x).
\end{equation}
Due to the translational invariance of the medium in the $y-$coordinate, 
if the source plane is translated to  $y=y_1$, then the spatial domain solution
${\bmi  b }(z,x,y; \omega)$ is translated to  ${\bmi  b }(z,x,y-y_1; 
\omega)$. Consequently, it is straight forward to derive the solution 
appropriate to general source conditions, by translating the field 
for a point source in the $y=0$ plane, accordingly.

In (\ref{govEqn0}) both $\omega$ and ${  k}_{  y}$ are {\em passive 
parameters\/}. The space/time solution is constructed from a 
superposition of solutions with different values of $\omega$ and ${  k}_{  y}$
as prescribed by the {\em inverse Fourier transforms\/} 
(\ref{InvLaplace}) on frequency and 
(\ref{kysyn}) on $k_y$. The main task remaining, is the construction the 
transformed displacement/traction vector ${\bmi  b }(z,x,k_y; \omega)$ in (\ref{govEqn0}). 
Often, we will not indicate the explicit dependence of  ${\bmi  b }(z,x,k_y; \omega)$ 
on $k_y$ and $\omega$ but write ${\bmi  b }(z,x)$ instead.

The representation (\ref{govEqn0}) is most convenient when the
dominant material variation is in the $z$ coordinate direction, as we have
singled out this coordinate for special treatment, allowing easy application of
boundary and interface conditions on fixed depth planes.  Also, the
displacement ${\bmi  u}$ and depth plane traction ${\bmi  t}$ are accorded equal status.
In the more traditional displacement formulation, the  offset $x$ and depth $z$
coordinates appear within the equations in a symmetric form (due to isotropy)
whereas the traction field must be derived from the displacement via the
elastic analogue of {\em Hooke's law\/} which involves further spatial
differentiations of $\bmi  u$.


%%%%%%%%%%%%%%%%%%%%%%%%%%%%%%%%%%%%%%%%%%%%%%%%%%%%%%%%%%%%%%%%%%%%%%%%%%%%%%%%%%%%%
%%%%%%%%%%%%%%%%%%%%%%%%%%%%%%%%%%%%%%%%%%%%%%%%%%%%%%%%%%%%%%%%%%%%%%%%%%%%%%%%%%%%%
We note in passing that (\ref{govEqn0}) simplifies considerably in the 
special case where the motion is driven by a {\em line source\/}, along 
the $y-$axis so that we may set  $k_y\equiv 0$ in (\ref{Adef}) which then 
reduces to
\[
{\CB A}(z,x;{\partial }_{x},0)
=\left[{\begin{array}{cccccc}
0&- \gamma {\partial }_{  x}&  0&  a&0&0\\ -{ \partial 
}_{x}&0&0&0&b&0\\ 0&0&0&0&0&b\\ - \rho {\omega }^{  
2}&  0&0&0&-{ \partial }_{x}&0\\ 0&- \rho {\omega 
}^{  2}  -{ \partial }_{x} \zeta {\partial }_{  
x} &  
0&  -{ \partial }_{x} \gamma &  0&0\\ 
0&0 &  - \rho {\omega }^{  2}  
 -{ \partial }_{x} \mu 
{\partial }_{  x}& 0 &  0&0\end{array}}\right].
\] 
In this instance,  (\ref{govEqn0}) may be decomposed into its {\em in plane\/} 
equation
\begin{equation} 	\label{PSVeqn}
{\frac{ \partial }{\partial z}}\left[{\begin{array}{c}{u}_{z}\\
{u}_{x}\\
{\tau }_{zz}\\
{\tau }_{xz}\end{array}}\right]
=\left[{\begin{array}{cccc}
0&- \gamma {\partial }_{  x}&  a&0\\ 
-{ \partial }_{x}&0&0&b\\
- \rho {\omega }^{  2}&  0&0&-{ \partial }_{x}\\ 
0&- \rho {\omega }^{  2}  -{ \partial }_{x} \zeta {\partial }_{  x} & 
 -{ \partial }_{x} \gamma &  0 
\end{array}}\right]
\left[{\begin{array}{c}{u}_{z}\\
{u}_{x}\\
{\tau }_{zz}\\
{\tau }_{xz}\end{array}}\right]
\end{equation}
for the $P/SV$ field and {\em out of plane\/} equation 
\begin{equation} 	\label{SHeqn}
	{\frac{\partial }{\partial z}}\left[{\begin{array}{c}{u}_{y}\\
	{\rm \tau }_{yz}\end{array}}\right]
	=\left[{\begin{array}{cc}
	0&b\\ 
	- \rho {\omega }^{  2}-{ \partial }_{x} \mu {\partial }_{  x}& 0
	\end{array}}\right]
	\left[{\begin{array}{c}{u}_{y}\\
	{\rm \tau }_{yz}\end{array}}\right]
\end{equation}
for the $SH$ field.
For the most part, we will concentrate on the solution of (\ref{govEqn0}) 
because the source conditions needed for the applicability (\ref{PSVeqn}) and (\ref{SHeqn})
are too restrictive.
%%%%%%%%%%%%%%%%%%%%%%%%%%%%%%%%%%%%%%%%%%%%%%%%%%%%%%%%%%%%%%%%%%%%%%%%%%%%%%%%%%%%%
%%%%%%%%%%%%%%%%%%%%%%%%%%%%%%%%%%%%%%%%%%%%%%%%%%%%%%%%%%%%%%%%%%%%%%%%%%%%%%%%%%%%%




Consistent with the decomposition of the material properties in (\ref{mdecomp}), we
 {\em split\/} (\ref{Adef}) into
\begin{equation} \label{Asplit}
{\CB A}({\bmi x};{\partial }_{x},{\partial }_{y};k_y)=
{\rf{\bmi A}}({\bmi x};{\partial }_{x},{\partial }_{y})
+{{\mit \Delta} {\bmi A}}({\bmi x};{\partial }_{x},{\partial }_{y}) 
\end{equation}
where
\begin{eqnarray} \label{Adefref}
\lefteqn{
{\rf{\bmi A}}(z,x;{\partial }_{x},{k}_{y})
=}\nonumber\\
&&\left[{\begin{array}{cccccc}
0&- {\rf \gamma} {\partial }_{  x}&  -i {\rf \gamma} {  k}_{  y}&  {\rf a}&0&0\\ -{ \partial 
}_{x}&0&0&0&{\rf b}&0\\ -i{k}_{y}&0&0&0&0&{\rf b}\\ - {\rf \rho} {\omega }^{  
2}&  0&0&0&-{ \partial }_{x}&-i{k}_{y}\\ 0&- {\rf \rho} {\omega 
}^{  2}  -{ \partial }_{x} {\rf \zeta} {\partial }_{  
x}  + {\rf \mu} {  k}_{  y}^{  2}&  
-i{k}_{y}{ \partial }_{x} {\rf \chi}   -i{k}_{y} {\rf \mu} {\partial 
}_{  x}&  -{ \partial }_{x} {\rf \gamma} &  0&0\\ 
0&-i{k}_{y} {\rf \chi} {\partial }_{  x}  -i{k}_{y}{ \partial 
}_{x} {\rf \mu} &  - {\rf \rho} {\omega }^{  2}  + {\rf \zeta }
{  k}_{  y}^{  2}  -{ \partial }_{x} {\rf \mu} 
{\partial }_{  x}&  -i{k}_{y} {\rf \gamma} &  
0&0\end{array}}\right]
\end{eqnarray}
where ${\rf \gamma}$ is the {\em reference medium\/} value of ${\gamma}$
from Table \ref{ATable} and similar definitions hold for the remaining 
{\em reference\/} material parameters 
in (\ref{Adefref}). Also,
\begin{eqnarray} \label{dAdef}
\lefteqn{
{{\mit \Delta} {\bmi A}}(z,x;{\partial }_{x},{k}_{y})
=}\\
&&\!\!\!\!\!\!\!\!\!\!\!\!
\left[{\begin{array}{cccccc}
0&- {{\mit \Delta} \gamma} {\partial }_{  x}&  -i {{\mit \Delta} \gamma} {  k}_{  y}&  
{{\mit \Delta} a}&0&0\\0&0&0&0&{{\mit \Delta} b}&0\\ 0&0&0&0&0&{{\mit \Delta} b}\\ - {{\mit \Delta} \rho} {\omega }^{  
2}&  0&0&0&0&0\\ 0&- {{\mit \Delta} \rho} {\omega 
}^{  2}  -{ \partial }_{x} {{\mit \Delta} \zeta} {\partial }_{  
x}  + {{\mit \Delta} \mu} {  k}_{  y}^{  2}&  
-i{k}_{y}{ \partial }_{x} {{\mit \Delta} \chi}   -i{k}_{y} {{\mit \Delta} \mu} {\partial 
}_{  x}&  -{ \partial }_{x} {{\mit \Delta} \gamma} &  0&0\\ 
0&-i{k}_{y} {{\mit \Delta} \chi} {\partial }_{  x}  -i{k}_{y}{ \partial 
}_{x} {{\mit \Delta} \mu} &  - {{\mit \Delta} \rho} {\omega }^{  2}  + {{\mit \Delta} \zeta }
{  k}_{  y}^{  2}  -{ \partial }_{x} {{\mit \Delta} \mu} 
{\partial }_{  x}&  -i{k}_{y} {{\mit \Delta} \gamma} &  0&0\end{array}}\right]\nonumber
\end{eqnarray}
with the material contrast parameters in (\ref{dAdef}) defined in Table \ref{contrastTable}.

\begin{table}
\begin{tabular*}{100mm}{ll}\hline
contrast & definition\\ \hline
 $ \Delta \rho$ & $\rho - {\rf \rho}$\\
 $ \Delta \mu$ & $\mu - {\rf \mu}$\\
 $ \Delta \gamma$ & $2{\frac{{ \rf \beta }{}^{2}}{{ \rf \alpha }{}^{2}}}-
 2{\frac{{\beta}^{2}}{{\alpha}^{2}}}$\\
$\Delta \zeta$ &$4 \rho{\beta 
}^{2}\left[{1-{\frac{{\beta  }^{2}}{{\alpha}^{2}}}}\right]-4{\rf \rho} 
{\rf \beta }{}^{  2}\left[{ 
 1-{\frac{{ \rf \beta }{}^{2}}{{\rf  \alpha }{}^{2}}}}\right]$\\
$\Delta \chi$ &$2 \rho {\beta}^{2}\left[{1-2{\frac{{ \beta  
}^{2}}{{\alpha}^{2}}}}\right]-2{\rf \rho} {\rf \beta }{}^{  2}\left[{ 
 1-2{\frac{{\rf \beta }{}^{2}}{{\rf \alpha }{}^{2}}}}\right]$\\
$\Delta  a$ & ${\frac{1}{\rho{\alpha}^{2}}}-
 {\frac{1}{ {\rf \rho} {{\rf \alpha} }{}^{  2}}}$\\ 
$\Delta  b$ & ${\frac{1}{\rho{\beta}^{2}}}-
 {\frac{1}{ {\rf \rho} {{\rf \beta} }{}^{  2}}}$\\  
 \hline
\end{tabular*}
\caption{\label{contrastTable} { The material contrast parameters}}
\end{table} 

Using the decomposition of (\ref{Asplit}) in the governing equation (\ref{govEqn0}) gives
\begin{equation} \label{Leqn}
{\bmi L}(z,x;{\partial }_{z},{\partial }_{x},{k}_{y})\cdot{\bmi b}(z,x,k_y; \omega) = 
\tilde{\bmi f} + {\mit \Delta} {\bmi A}(z,x;{\partial }_{x},{k}_{y})\cdot{\bmi b}(z,x,k_y; \omega)
\end{equation}
where the {\em reference medium\/} operator ${\bmi L}$ is defined as
\begin{equation} \label{Lref} 
{\bmi L}(z,x;{\partial }_{z},{\partial }_{x},{k}_{y})\equiv
\left[\partial_{z}-\rf{\bmi A}(z,x;{\partial }_{x},{k}_{y})\right]
\end{equation}
with $\rf{\bmi A}$ given in (\ref{Adefref}).
Scrupulously indicating all the arguments of the various operators introduced in 
this thesis can be very cumbersome, so for example, we will often write
${\bmi L}$ to mean ${\bmi 
L}(z,x;{\partial }_{z},{\partial }_{x},{k}_{y})$ when no misunderstanding 
is likely.

\section{The $2-$D Elastic Representation Theorem} \label{section2Drep}
Given the linear differential operator ${\bmi L}$ of (\ref{Leqn}), it is possible to
derive an {\em adjoint} operator ${\bmi L}^{\dagger}$ relative to an 
arbitrary {\em bilinear form} $(\cdot,\cdot)$, so that (Fig.\ref{G2DThm}) 
\begin{eqnarray} \label{rep1} 
({\bmi b}_{2},{\bmi L}{\bmi b}_{1})-({\bmi L}^{\dagger}{\bmi b}_{2},{\bmi b}_{1})&=& 
\int \left( \partial_{z}{\bmi
\Phi}_{z}+\partial_{x}{\bmi \Phi}_{x}\right)  \,  { \rm d}z { \rm d}x\nonumber\\
&=&\oint \left( {\bmi \Phi}_{z} \ { \rm d}x - {\bmi \Phi}_{x}\  { \rm d}z \right)
\end{eqnarray} 
for any pair of sufficiently differentiable, $6-$component vector
fields ${\bmi b}_{1}(z,x)$ and ${\bmi b}_{2}(z,x)$, not necessarily solutions of
${\bmi L}$ or its adjoint ${\bmi L}^{\dagger}$ and ${\bmi \Phi}=({\bmi 
\Phi}_z,{\bmi \Phi}_x)^T$ 
is a yet to be determined $2-$component {\em flux\/} field. Here we adopt the specific bilinear
form
$(,)_{{\mit \Omega}}$, defined over the scattering region
${\mit \Omega}$ according to
\begin{equation} \label{bilin} 
({\bmi b}_{1},{\bmi b}_{2})_{{\mit \Omega}}\equiv
\int_{{\mit \Omega}} {\rm d}{\bmi x}  \left\langle{\bmi b}_{1}({\bmi x}), {\bmi 
b}_{2}({\bmi x})\right\rangle  
\end{equation}
where $\left\langle{\bmi b}_{1}({\bmi x}), {\bmi b}_{2}({\bmi x})\right\rangle$
is defined as 
\begin{equation}	\label{sharpdef}
\left\langle{\bmi b}_{1}({\bmi x}), {\bmi b}_{2}({\bmi x})\right\rangle\equiv
{\bmi b}_{1}^{T}{\bmi J}_{6}{\bmi b}_{2}({\bmi x})
\end{equation}
and ${\bmi J}_{6}$ is given by
\begin{equation}	\label{Jdef}
{\bmi J}_{6}=\left[{\begin{array}{cc}
{\bf 0}_{3}&{\bmi I}_{3}\\
-{\bmi I}_{3}&{\bf 0}_{3}\end{array}}\right].
\end{equation}

\begin{figure}
\HideDisplacementBoxes
%\ShowDisplacementBoxes
\centerline{\BoxedEPSF{G2DThm.EPSF scaled 800}}
\caption{The $2-$D Green's Theorem}
\protect\label{G2DThm}
\end{figure}

Using the {\em product rule for partial differentiation\/},
\[
{\frac{ \partial   \ a(x)\ b(x)}{ \partial   \, 
x}}\ =\ b(x)\ {\frac{ \partial   \ a(x)}{ \partial   
\, x}}\ +\ a(x)\ {\frac{ \partial   \ b(x)}{ \partial  
 \, x}} 
\] 
of differentiable functions $a(x)$ and $b(x)$, it can be verified that 
\begin{equation} \label{partialint3}
\left\langle{{\bmi b}_{1},{\bmi L}{\bmi b}_{2}}\right\rangle-\left\langle{ 
{\bmi L}^{\dagger} {\bmi b}_{1},
{\bmi b}_{2}}\right\rangle
={ \partial }_{z}\left\langle{\bmi b}_{1},{\bmi b}_{2}\right\rangle
+{ \partial }_{x}\left\langle{\bmi T}_x(-k_y){\bmi b}{}_{1},
{\bmi T}_x(k_y){\bmi b}{}_{2}\right\rangle
\end{equation} 
where ${\bmi L}^{\dagger}$ is the operator adjoint of ${\bmi L}$, defined as
\begin{equation} \label{Ladjdef}
{\bmi L}^{\dagger}({\bmi x};{\partial }_{z},{\partial }_{x},{k}_{y})\equiv 
-{\bmi L}({\bmi x};{\partial }_{z},{\partial }_{x},-{k}_{y})
\end{equation}
and
\begin{equation} \label{Tractdef}
{\bmi T}_x(k_y)=\left[{\begin{array}{cccccc}
1&0&0&0&0&0\\
0&1&0&0&0&0\\
0&0&1&0&0&0\\
0&0&0&0&1&0\\
0&{\rf \zeta} {\partial }_{ x}& i\,{k}_{y} {\rf \chi} &{\rf \gamma} & 0&0\\
0&i\,{k}_{y}{\rf \mu} &{\rf \mu} {\partial }_{ x}& 0&0&0\end{array}}\right].
\end{equation}



Integrating both sides of (\ref{partialint3}) over a simply connected
spatial region ${\mit \Omega}$, gives 
\begin{eqnarray} 	\label{int1}
\lefteqn{
\left({{\bmi b}_{1},{\bmi L}\cdot{\bmi b}_{2}}\right)_{\mit \Omega}-
\left({{\bmi L}^{\dagger} \cdot{\bmi b}_{1},
{\bmi b}_{2}}\right)_{\mit \Omega}}\nonumber\\
&&=
\int_{\mit \Omega}  \left[
\left\langle{{\bmi b}_{1},{\bmi L}({\bmi x};{\partial }_{z},{\partial 
}_{x},{k}_{y}) \cdot{\bmi b}_{2}}\right\rangle 
+\left\langle{{\bmi L}({\bmi x};{\partial }_{z},{\partial 
}_{x},-{k}_{y}) \cdot{\bmi b}_{1},{\bmi b}_{2}}\right\rangle \right]\,{\rm d}{\bmi x} \nonumber\\
	&&=\int_{\mit \Omega} \left[
	{ \partial}_{z}\left\langle{{\bmi b}_{1},{\bmi b}_{2}}\right\rangle
	+ { \partial}_{x}\left\langle{{\bmi T}_x(- k_y){\bmi b}_{1},{\bmi T}_x ( k_y){\bmi 
	b}_{2}}\right\rangle \right]\, {\rm d}{\bmi x}\nonumber\\
	&&= \oint_{\partial\Omega} {\rm d}{x}
	\left\langle{{\bmi b}_{1},{\bmi b}_{2}}\right\rangle -
	{\rm d}{z} \left\langle{{\bmi T}_x(- k_y){\bmi b}_{1},{\bmi T}_x(k_y){\bmi b}_{2}}\right\rangle
\end{eqnarray}
where the last equation in (\ref{int1}) follows from the $2-$D version of Green's theorem.
The integral expression (\ref{int1}) is a {\em balance law\/} for the 
two component {\em flux field\/}
\[
({\bmi \Phi}_z,{\bmi \Phi}_x) \equiv (\left\langle{{\bmi b}_{1},{\bmi b}_{2}}\right\rangle,
\left\langle{{\bmi T}_x(- k_y){\bmi b}_{1},{\bmi T}_x( k_y){\bmi b}_{2}}\right\rangle )
\]
generated by the {\em virtual  source density\/}
\begin{eqnarray}
\lefteqn{	
\left\langle{{\bmi b}_{1},{\bmi L}({\bmi x};{\partial }_{z},{\partial 
}_{x},{k}_{y}) \cdot{\bmi 
	b}_{2}}\right\rangle
	+\left\langle{{\bmi L}({\bmi x};{\partial }_{z},{\partial 
}_{x},-{k}_{y}) \cdot{\bmi b}_{1},{\bmi 
b}_{2}}\right\rangle\equiv}\nonumber\\
&&	
\overbrace{\left\langle{{\bmi b}_{1},{\bmi L}({\bmi x};{\partial }_{z},{\partial 
}_{x},{k}_{y}) \cdot{\bmi 
	b}_{2}}\right\rangle}^{\mbox{source density}}-
\overbrace{\left\langle{{\bmi L}^{\dagger}({\bmi x};{\partial }_{z},{\partial 
}_{x},{k}_{y}) \cdot{\bmi b}_{1},{\bmi 
b}_{2}}\right\rangle}^{\mbox{sink density}}
\end{eqnarray}
({\em see\/} Fig.\ref{G2DThm}).

Equation (\ref{int1}) will be used extensively in the development of this 
chapter, and to save space, we rewrite it in the following compact notation
\begin{equation} \label{rep2}
\fbox{$ \displaystyle
\left({\bmi b}_{1},{\bmi L}{\bmi b}_{2}\right)_{{\mit \Omega}}-\left({\bmi 
L}^{\dagger}{\bmi b}_{1},{\bmi b}_{2}\right)_{{\mit \Omega}}=
\left\langle{\bmi b}_{1},{\bmi b}_{2}
\right\rangle_{\partial\Omega}
$} 
\end{equation} 
where the boundary integral $\langle{\bmi b}_{1},{\bmi b}_{2}\rangle_{\partial\Omega}$
 is given by the  cartesian form
\begin{eqnarray} \label{Bi}
\lefteqn{\langle{\bmi b}_{1},{\bmi b}_{2}\rangle_{\partial\Omega} =
\oint_{\partial {\mit \Omega}} \left[ {\bmi b}_{1}^{T}(z,x) {\bmi J}_{6} {\bmi
b}_{2}(z,x) \  {\rm d}x - ({\bmi T}_x(-k_y){\bmi b}_{1})^{T}(z,x) {\bmi J}_{6} 
{\bmi T}_x(k_y){\bmi b}_{2}(z,x) \ {\rm d}z \right]}\nonumber\\
&&
\end{eqnarray}
or the {\em parametric} form
\begin{equation} \label{Bi2}
\langle{\bmi b}_{1},{\bmi b}_{2}\rangle_{\partial\Omega} =
\oint_{\partial {\mit \Omega}} 
{{\bmi b}^{T}_{1}(z(s),x(s);-k_y,\hat{\bmi n})}  {\bmi J}_{6} 
{{\bmi b}_{2}(z(s),x(s);k_y,\hat{\bmi n})} { \rm d}s
\end{equation}
where ${ \rm d}s$ is the  arc length of the boundary $\partial {\mit \Omega}$, 
\[
\hat{\bmi n} \equiv {\bmi e}_{\mit z}{\frac{{\rm d\,}x}{{\rm d\,}s}}
 -{\bmi e}_{x}{\frac{{\rm d\,}z}{{\rm d\,}s}}
\] 
is the outward unit normal (Fig.\ref{G2DThm}) and
\begin{equation} \label{normalBdef}
{\bmi b}_{1,2}(z,x;\hat{\bmi n})={\bmi N}({\hat{\bmi n}}){\bmi b}_{1,2}(z,x),
\end{equation}
where
\begin{equation} \label{ndefMS}
{\bmi N}(k_y,{\hat{\bmi n}})=
\left[{\begin{array}{cccccc}1&0&0&0&0&0\\
0&1&0&0&0&0\\
0&0&1&0&0&0\\
0&0&0&{\frac{d\,x}{d\,s}}&-{\frac{d\,z}{d\,s}}&0\\
0&-{\frac{d\,z}{d\,s}} {\rf \zeta} {\partial }_{  x}&  -i{k}_{y}\,{\frac{d\,z}{d\,s}} 
{\rf \chi} &  -{\frac{d\,z}{d\,s}} {\rf \gamma} &{\frac{  d\,x}{d\,s}}&  0\\
0&-i{k}_{y}\,{\frac{d\,z}{d\,s}} {\rf \mu} &  -{\frac{d\,z}{d\,s}} {\rf \mu} 
{\partial }_{  x}&  
0&0&{\frac{d\,x}{d\,s}}\end{array}}\right].
\end{equation}



\section{Reciprocity}
The great generality of (\ref{rep2}) derives from its abstract nature, as 
the {\em test } fields ${\bmi b}_{1}$ and ${\bmi b}_{2}$ in
(\ref{rep2}) have been
arbitrary, subject only to sufficient differentiability to ensure the validity
of the various integration by parts operations needed to derive (\ref{rep2}).
We now impose further conditions on these fields. Specifically let us 
assume they are point source  solutions in a halfspace with no 
inclusions, and continuous elastic properties, so that
\begin{equation} \label{prob0}
\left.{\begin{array}{lr}
{\bmi L}^{\dagger}\cdot
{\rf{\bmi  G}}_{1}({\bmi x},{\bmi x}_{1},k_y) = -{\bmi I}_{6}\delta({\bmi 
x}-{\bmi x}_{1}),&
{\bmi x}_{1}\in{ R}^{2}\\
{\bmi L}\cdot
{\rf{\bmi  G}}_{2}({\bmi x},{\bmi x}_{2},k_y) = {\bmi I}_{6}\delta({\bmi 
x}-{\bmi x}_{2}),&{\bmi x}_{2}\in{ R}^{2}
\end{array}}\right\},
\end{equation}
subject to the traction free condition at the free surface, and the 
radiation condition at infinity. The differential operator ${\bmi L}$ was 
given in (\ref{Lref}) and the 
corresponding adjoint operator  ${\bmi L}^{\dagger}$ was given in 
(\ref{Ladjdef}).
Note that the {\em point source solutions\/} ${\rf{\bmi  G}}_{1}({\bmi 
x},{\bmi x}_{1},k_y)$ and ${\rf{\bmi  G}}_{2}({\bmi x},{\bmi x}_{2},k_y)$ both have 
two arguments. The right most argument is the {\em source\/} coordinate, and the 
left argument is the {\em field\/} coordinate, so that ${\rf{\bmi  G}}_{1}({\bmi 
x},{\bmi x}_{1},k_y) \left( {\rf{\bmi  G}}_{2}({\bmi 
x},{\bmi x}_{2},k_y) \right)$ has a source at ${\bmi x}_1 ( {\bmi x}_2)$ and 
a field coordinate at ${\bmi x}$. 

Then (\ref{rep2}) gives
\begin{eqnarray} \label{rep0}
\lefteqn{
\left({\rf{\bmi  G}}_{1}({\bmi x},{\bmi x}_{1},k_y),{\bmi L}{\rf{\bmi  G}}_{2}({\bmi x},{\bmi 
x}_{2},k_y)\right)_{{\mit \Omega}}-\left({\bmi
L}^{\dagger}{\rf{\bmi  G}}_{1}({\bmi x},{\bmi x}_{1},k_y),{\rf{\bmi  G}}_{2}({\bmi x},{\bmi 
x}_{2},k_y)\right)_{{\mit \Omega}}}\nonumber\\
&&
\equiv\left\langle{\rf{\bmi  G}}_{1}({\bmi x},{\bmi x}_{1},k_y),{\rf{\bmi  
G}}_{2}({\bmi x},{\bmi x}_{2},k_y)
\right\rangle_{\partial{{\mit \Omega}}^{\scriptscriptstyle \infty + f}}
\equiv{\bf 0}
\end{eqnarray}
where the boundary integral 
$
\left\langle{\rf{\bmi  G}}_{1}({\bmi x},{\bmi x}_{1},k_y),
{\rf{\bmi  G}}_{2}({\bmi x},{\bmi x}_{2},k_y)\right\rangle_{\partial{{\mit 
\Omega}}^{\infty+f}}
$
 is over a surface
$\partial{{\mit \Omega}}^{{\scriptscriptstyle \infty}}$ at infinity and over the free surface 
$\partial{{\mit \Omega}}^{f}$ and consequently, it vanishes in 
accordance with the {\em radiation condition} and the {\em free surface 
condition\/}. Indeed the annihilation of
\[
\left\langle{\rf{\bmi  G}}_{1}({\bmi x},{\bmi x}_{1},k_y),{\rf{\bmi  G}}_{2}({\bmi x},{\bmi 
x}_{2},k_y) \right\rangle_{\partial{{\mit \Omega}}^{{\scriptscriptstyle \infty}}}
\]
may be taken as the definition of the radiation condition.
Using (\ref{prob0}) the left hand side of (\ref{rep0}) is
\begin{eqnarray}	\label{recip0}
\lefteqn{\left({\rf{\bmi  G}}_{1}({\bmi x},{\bmi x}_{1},k_y),{\bmi I}_{6}\delta({\bmi 
	x}-{\bmi x}_{2})\right)_{{\mit \Omega}}+\left({\bmi I}_{6}\delta({\bmi 
	x}-{\bmi x}_{1}),{\rf{\bmi  G}}_{2}({\bmi x},{\bmi x}_{2},k_y)\right)_{{\mit 
	\Omega}}=}\nonumber\\
&&
\rf{\bmi  G}{}_{1}^{T}({\bmi x}_{2};{\bmi x}_{1},k_y){\bmi J}_{6} +
{\bmi J}_{6} {\rf{\bmi  G}}_{2}({\bmi x}_{1};{\bmi x}_{2},k_y) \equiv{\bf 0}
\end{eqnarray}
or upon rearranging we get
\begin{equation}	\label{recipDef0}
{\rf{\bmi  G}}_{1}({\bmi x}_{2};{\bmi x}_{1},k_y) \equiv
{\bmi J}_{6}\cdot\rf{\bmi  G}{}^{T}_{2}({\bmi x}_{1};{\bmi x}_{2},k_y)\cdot{\bmi J}_{6}
\end{equation}
From the governing equations (\ref{prob0}) and the definition of the adjoint 
operator ${\bmi L}^{\dagger}$ (\ref{Ladjdef})
\[
{\rf{\bmi  G}}_{1}({\bmi x}_{2};{\bmi x}_{1}; k_y)\equiv 
{\rf{\bmi  G}}_{2}({\bmi x}_{2};{\bmi x}_{1}; -k_y)
\]
so that
\begin{equation}	\label{recipDef}
\fbox{$ \displaystyle
{\rf{\bmi  G}}_2({\bmi x}_{2};{\bmi x}_{1};k_y) \equiv
{\bmi J}_{6}\cdot\rf{\bmi  G}{}^{T}_2({\bmi x}_{1};{\bmi x}_{2};-k_y)\cdot{\bmi 
J}_{6}
$}\end{equation}
which is the {\em reciprocity} relation for point source solutions in 
free space.


Very often, we split the Green's tensor ${\rf{\bmi  G}}({\bmi 
x}_{2};{\bmi x}_{1}; k_y)$ into the product 
\[
{\rf{\bmi  G}}({\bmi 
x}_{2};{\bmi x}_{1}; k_y)\equiv {\rf{\bmi  P}}({\bmi 
x}_{2};{\bmi x}_{1}; k_y){\bmi J}_6
\]
where ${\rf{\bmi  P}}({\bmi 
x}_{2};{\bmi x}_{1}; k_y)$ is defined entirely by its relation to
${\rf{\bmi  G}}({\bmi 
x}_{2};{\bmi x}_{1}; k_y)$. We refer to ${\rf{\bmi  P}}({\bmi 
x}_{2};{\bmi x}_{1}; k_y)$ as the {\em propagator\/}. It follows directly 
from (\ref{recipDef}) that  the propagator 
obeys the simple reciprocity relation
\begin{equation}	\label{recipPDef}
\fbox{$ \displaystyle
{\rf{\bmi  P}}({\bmi x}_{2};{\bmi x}_{1};k_y) \equiv
\rf{\bmi  P}\!{}^{T}({\bmi x}_{1};{\bmi x}_{2};-k_y).
$}\end{equation}



%%%%%%%%%%%%%%%%%%%%%%%%%%%%%%%%%%%%%%%%%%%%%%%%%%%%%%%%%%%%%%%%%%%%%%%%%%%%%%%%%

In the sections that follow, we encounter numerous instances of
\[
{\bmi J}_{6}\left[{\bmi T}_{ x}(-k_y)
\rf{\bmi  G}({\bmi x};{\bmi x}^{\prime};-k_y)\right]^{T}{\bmi J}_{6}
\]
for various (receiver, source) combinations $({\bmi x};{\bmi x}^{\prime})$, 
where the operator ${\bmi T}_{ x}$ was given in (\ref{Tractdef}) and acts on the 
field (or receiver) coordinate ${\bmi x}$ in $\rf{\bmi  G}({\bmi x};{\bmi 
x}^{\prime};-k_y)$. It follows from (\ref{recipDef}) that:
\begin{equation}	\label{recipTDef}
\fbox{$ \displaystyle
\stackrel{\scriptscriptstyle _s}{\bmi T}{}_{\!\!\! x}(-k_y)
\rf{\bmi  G}{}({\bmi x}_{2};{\bmi x}_{1};k_y)\equiv
{\bmi J}_{6}\left[\stackrel{\scriptscriptstyle _r}{\bmi T}{}_{\!\!\! 
x}(-k_y)
\rf{\bmi  G}{}({\bmi x}_{1};{\bmi x}_{2};-k_y)\right]^{T}{\bmi J}_{6}
$}\end{equation}
where $\stackrel{\scriptscriptstyle _s}{\bmi T}{}_{\!\!\! x}$ operates on 
the source vector ${\bmi x}_{1}$ vector in $\rf{\bmi  G}({\bmi x}_{2};{\bmi x}_{1})$, and 
we have written $\stackrel{\scriptscriptstyle _r}{\bmi T}{}_{\!\!\! x}
\rf{\bmi  G}({\bmi x}_{1};{\bmi x}_{2})$ in place of the previously 
employed notation ${\bmi T}_{ x}\rf{\bmi  G}({\bmi x}_{1};{\bmi 
x}_{2})$ to indicate an operation on the field (or receiver) coordinate 
${\bmi x}_{1}$.


\section{The Representation Theorems}
\subsection{Exterior Representation Theorem}
{\noindent{\bf Theorem 1}} Given that
 \begin{eqnarray*}
\lefteqn{{\bmi L}({\bmi x},\partial_{z},\partial_{x};k_y) {\bmi G}({\bmi
x},{\bmi x}_{S};k_y) =}\nonumber \\
&&\qquad
{\bmi I}_{6}\delta({\bmi x}-{\bmi x}_{S}) + 
{\chi}_{{{\mit \Omega}}} {\mit \Delta} {\bmi A}({\bmi
x},\partial_{x};k_y) {\bmi G}({\bmi x},{\bmi x}_{S};k_y),\qquad {\bmi 
x}_{S}\in{\mit \Omega}^{e}
\end{eqnarray*}
and
\begin{displaymath}
{\bmi L}^{\dagger}({\bmi x},\partial_{z},\partial_{x};k_y) \rf{\bmi G}({\bmi
x},{\bmi x}_{R};-k_y) = -{\bmi I}_{6}\delta({\bmi x}-{\bmi x}_{R}),\ \ \ \ {\bmi 
x}_{R}\in{R}^{2}
\end{displaymath}
then
\begin{eqnarray*}
\lefteqn{
\left.
\begin{array}{cc}
{{\bmi G}}({\bmi x}_{R};{\bmi x}_{S};k_y)& {\bmi x}_{R}\in {\mit \Omega}^{e}\\
{\bf 0}&{\bmi x}_{R}\in {\mit \Omega}
\end{array}  \right\}
={\rf{\bmi G}}({\bmi x}_{R};{\bmi x}_{S};k_y)
}\\
&&
\qquad\qquad\qquad\qquad\qquad\qquad
-\oint_{\partial {\mit \Omega}^{+}} \left[ {\rf{\bmi G}}({\bmi x}_{R};{z, x};k_y) 
{{\bmi G}}({z, x};{\bmi x}_{S};k_y) \  {\rm d}x \right.\\
&&
\qquad\qquad\qquad\qquad\qquad\qquad
- \left.\stackrel{\scriptscriptstyle _s}{\bmi T}{}_{\!\!\! x}(-k_y){\rf{\bmi 
G}}({\bmi x}_{R};{z, x};k_y)  
\stackrel{\scriptscriptstyle _r}{\bmi T}{}_{\!\!\! x}(k_y){{\bmi G}}({z, 
x};{\bmi x}_{S};k_y) \ {\rm d}z \right]
\end{eqnarray*}

{\em Proof}
It follows from (\ref{rep2}) with ${\mit \Omega}$ chosen as ${\mit 
\Omega^{e}}$ in that 
equation, that
\begin{eqnarray} \label{rep02}
\lefteqn{
\left({\rf{\bmi G}}({\bmi x};{\bmi x}_{R};-k_y),{\bmi L}\cdot{{\bmi 
G}}({\bmi x};{\bmi x}_{S};k_y)\right)_{{\mit \Omega}^{e}}-
\left({\bmi L}^{\dagger}\cdot{\rf{\bmi G}}({\bmi x};{\bmi x}_{1};-k_y),{ 
{\bmi G}}({\bmi x};{\bmi x}_{2};k_y)\right)_{{\mit \Omega}^{e}}=}\nonumber\\
&&
\left\langle{\rf{\bmi G}}({\bmi x};{\bmi x}_{1};-k_y),{{\bmi 
G}}({\bmi x};{\bmi x}_{2};k_y)\right\rangle_{\partial{\mit \Omega}^{+}}
+\overbrace{\left\langle{\rf{\bmi G}}({\bmi x};{\bmi x}_{1};-k_y),
{{\bmi G}}({\bmi x};{\bmi 
x}_{2};k_y)\right\rangle_{\partial{\mit \Omega}^{{\scriptscriptstyle \infty + 
f}}}}^{\equiv {\bmi 0}\ \small
\mbox{free-surface/radiation  conditions}}.
\end{eqnarray} 


We can write (\ref{rep02}) more explicitly as
\begin{eqnarray} \label{rep02explicit}
\lefteqn{
\int_{{\mit \Omega}^{e}} {\rf{\bmi G}}{}^{T} ({\bmi x};{\bmi 
x}_{R};-k_y){\bmi J}_{6}{\bmi L}\cdot{{\bmi G}}({\bmi 
x};{\bmi x}_{S};k_y) {\rm d}{\bmi x} -
}\nonumber\\
&&
\int_{{\mit \Omega}^{e}} ({\bmi L}^{\dagger}\cdot{\rf{\bmi G}})^{T}({\bmi x};
{\bmi x}_{R};-k_y){\bmi J}_{6}{{\bmi G}}({\bmi 
x};{\bmi x}_{S};k_y) {\rm d}{\bmi x}=\nonumber\\
&&
\oint_{\partial {\mit \Omega}^{+}} \left[ {\rf{\bmi G}}{}^{T}({z, x};{\bmi 
x}_{R};-k_y) {\bmi J}_{6} 
{{\bmi G}}({z, x};{\bmi x}_{S};k_y) \  {\rm d}x\right.\nonumber\\
&&\left. 
- ({\bmi T}_x(-k_y){\rf{\bmi G}}({z, x};{\bmi x}_{R};-k_y))^{T} {\bmi J}_{6} 
{\bmi T}_x(k_y){{\bmi G}}({z, x};{\bmi x}_{S};k_y) \ {\rm d}z \right]
\end{eqnarray} 
But multiplying (\ref{rep02explicit}) through by ${\bmi J}_{6}$ and using 
the reciprocity relations (\ref{recipDef}) and (\ref{recipTDef}) we get
\begin{eqnarray} \label{rep2explicit}
\lefteqn{
\int_{{\mit \Omega}^{e}} {\rf{\bmi G}}({\bmi x}_{R};{\bmi x};k_y){\bmi L}{{\bmi G}}({\bmi 
x};{\bmi x}_{S};k_y) {\rm d}{\bmi x} -
\int_{{\mit \Omega}^{e}} {\bmi J}_{6} ({\bmi L}^{\dagger}{\rf{\bmi G}})^{T}(
{\bmi x};{\bmi x}_{R};-k_y){\bmi J}_{6}{{\bmi G}}({\bmi 
x};{\bmi x}_{S};k_y) {\rm d}{\bmi x}}\nonumber\\
&&=
\oint_{\partial {\mit \Omega}^{+}} \left[ {\rf{\bmi G}}({\bmi x}_{R};{z, x};k_y) 
{{\bmi G}}({z, x};{\bmi x}_{S};k_y) \  {\rm d}x \right.\nonumber\\
&&\qquad\qquad\qquad\qquad
\left.
- \stackrel{\scriptscriptstyle _s}{\bmi T}{}_{\!\!\! x}(-k_y){\rf{\bmi 
G}}({\bmi x}_{R};{z, x};k_y)  
\stackrel{\scriptscriptstyle _r}{\bmi T}{}_{\!\!\! x}(k_y){{\bmi G}}({z, 
x};{\bmi x}_{S};k_y) \ {\rm d}z \right]
\end{eqnarray}
Using 
\[
\int_{{\mit \Omega}^{e}} \overbrace{{\bmi J}_{6} ({\bmi L}^{\dagger}{\rf{\bmi G}})^{T}(
{\bmi x};{\bmi x}_{R};-k_y){\bmi J}_{6}}^{{\bmi I}_{6}\delta({\bmi x}-{\bmi 
x}_R)}{{\bmi G}}({\bmi 
x};{\bmi x}_{S};k_y) {\rm d}{\bmi x}=\left\{
\begin{array}{cc}
{{\bmi G}}({\bmi x}_{R};{\bmi x}_{S};k_y)& {\bmi x}_{R}\in {\mit \Omega}^{e}\\
{\bf 0}&{\bmi x}_{R}\in {\mit \Omega}
\end{array}  \right.
\]
and
\[
\int_{{\mit \Omega}^{e}} {\rf{\bmi G}}({\bmi x}_{R};{\bmi x};k_y)
\overbrace{{\bmi L}{{\bmi G}}({\bmi x};{\bmi x}_{S};k_y)}^{{\bmi I}_{6}\delta({\bmi x}-{\bmi 
x}_S)} {\rm d}{\bmi x}=
{\rf{\bmi G}}({\bmi x}_{R};{\bmi x}_{S};k_y)
\]
in (\ref{rep2explicit}) and rearranging proves Thm.1.
\hfill \rule{11 pt}{11 pt}


%%%%%%%%%%%%%%%%%%%%%%%%%%%%%%%%%%%%%%%%%%%%%%%%%%%%%%%%%%%%%%%%%%%%%%%%%%%%%%%%%%%%%%%%%
\subsection{Interior Representation Theorem}
{\noindent{\bf Theorem 2}} Given that
 \begin{eqnarray*}
\lefteqn{{\bmi L}({\bmi x},\partial_{z},\partial_{x};k_y) {\bmi G}({\bmi
x},{\bmi x}_{S};k_y) =}\\
&&\qquad {\bmi I}_{6}\delta({\bmi x}-{\bmi x}_{S}) + 
{\chi}_{{{\mit \Omega}}} {\mit \Delta} {\bmi A}({\bmi
x},\partial_{x};k_y) {\bmi G}({\bmi x},{\bmi x}_{S};k_y),\qquad {\bmi 
x}_{S}\in{\mit \Omega}^{e}
\end{eqnarray*}
and
\begin{displaymath}
{\bmi L}^{\dagger}({\bmi x},\partial_{z},\partial_{x};k_y) \rf{\bmi G}({\bmi
x},{\bmi x}_{R};-k_y) = -{\bmi I}_{6}\delta({\bmi x}-{\bmi x}_{R}),\ \ \ \ {\bmi 
x}_{R}\in{R}^{2}
\end{displaymath}
then

\begin{samepage}
\begin{eqnarray*}
\lefteqn{
\left.
\begin{array}{cc}
{{\bmi G}}({\bmi x}_{R};{\bmi x}_{S};k_y)& {\bmi x}_{R}\in {\mit \Omega}\\
{\bf 0}&{\bmi x}_{R}\in {\mit \Omega}^{e}
\end{array}  \right\}
=\int_{{\mit \Omega}} {\rf{\bmi G}}({\bmi x}_{R};{\bmi x};k_y)
{\mit \Delta} {\bmi A}({\bmi
x},\partial_{x};k_y) {\bmi G}({\bmi x},{\bmi x}_{S};k_y) {\rm d}{\bmi x}
}\\
&&
\qquad\qquad\qquad\qquad\qquad\qquad
-\oint_{\partial {\mit \Omega}^{-}} \left[ {\rf{\bmi G}}({\bmi x}_{R};{z, x};k_y) 
{{\bmi G}}({z, x};{\bmi x}_{S};k_y) \  {\rm d}x \right.\\
&&
\qquad\qquad\qquad\qquad\qquad\qquad
- \left.\stackrel{\scriptscriptstyle _s}{\bmi T}{}_{\!\!\! x}(-k_y){\rf{\bmi 
G}}({\bmi x}_{R};{z, x};k_y)  
\stackrel{\scriptscriptstyle _r}{\bmi T}{}_{\!\!\! x}(k_y){{\bmi G}}({z, 
x};{\bmi x}_{S};k_y) \ {\rm d}z \right]
\end{eqnarray*}
\end{samepage}

{\em Proof}
It follows from (\ref{rep2}) that 
\begin{eqnarray} \label{thm20}
\lefteqn{
\left({\rf{\bmi G}}({\bmi x};{\bmi x}_{R};-k_y),{\bmi L}\cdot{{\bmi 
G}}({\bmi x};{\bmi x}_{S};k_y)\right)_{{\mit \Omega}}-
\left({\bmi L}^{\dagger}\cdot{\rf{\bmi G}}({\bmi x};{\bmi x}_{1};-k_y),{ 
{\bmi G}}({\bmi x};{\bmi x}_{2};k_y)\right)_{{\mit \Omega}}=}\nonumber\\
&&
\left\langle{\rf{\bmi G}}({\bmi x};{\bmi x}_{1};-k_y),{{\bmi 
G}}({\bmi x};{\bmi x}_{2};k_y)\right\rangle_{\partial{\mit \Omega}^{-}}.
\end{eqnarray} 
We can write (\ref{thm20}) more explicitly as
\begin{eqnarray} \label{rep20explicit}
\lefteqn{
\int_{{\mit \Omega}} {\rf{\bmi G}}{}^{T} ({\bmi x};{\bmi 
x}_{R};-k_y){\bmi J}_{6}{\bmi L}\cdot{{\bmi G}}({\bmi 
x};{\bmi x}_{S};k_y) {\rm d}{\bmi x} -
}\nonumber\\
&&
\int_{{\mit \Omega}} ({\bmi L}^{\dagger}\cdot{\rf{\bmi G}})^{T}({\bmi x};
{\bmi x}_{R};-k_y){\bmi J}_{6}{{\bmi G}}({\bmi 
x};{\bmi x}_{S};k_y) {\rm d}{\bmi x}=\nonumber\\
&&
\oint_{\partial {\mit \Omega}^{-}} \left[ {\rf{\bmi G}}{}^{T}({z, x};{\bmi 
x}_{R};-k_y) {\bmi J}_{6} 
{{\bmi G}}({z, x};{\bmi x}_{S};k_y) \  {\rm d}x\right.\nonumber\\
&&\left. 
- ({\bmi T}_x(-k_y){\rf{\bmi G}}({z, x};{\bmi x}_{R};-k_y))^{T} {\bmi J}_{6} 
{\bmi T}_x(k_y){{\bmi G}}({z, x};{\bmi x}_{S};k_y) \ {\rm d}z \right]
\end{eqnarray} 
But multiplying (\ref{thm20}) through by ${\bmi J}_{6}$ and using 
the reciprocity relations (\ref{recipDef}) and (\ref{recipTDef}) we get
\begin{eqnarray} \label{thm21}
\lefteqn{
\int_{{\mit \Omega}} {\rf{\bmi G}}({\bmi x}_{R};{\bmi x};k_y){\bmi L}{{\bmi G}}({\bmi 
x};{\bmi x}_{S};k_y) {\rm d}{\bmi x} -
\int_{{\mit \Omega}} {\bmi J}_{6} ({\bmi L}^{\dagger}{\rf{\bmi G}})^{T}(
{\bmi x};{\bmi x}_{R};-k_y){\bmi J}_{6}{{\bmi G}}({\bmi 
x};{\bmi x}_{S};k_y) {\rm d}{\bmi x}}\nonumber\\
&&=
\oint_{\partial {\mit \Omega}^{-}} \left[ {\rf{\bmi G}}({\bmi x}_{R};{z, x};k_y) 
{{\bmi G}}({z, x};{\bmi x}_{S};k_y) \  {\rm d}x \right.\nonumber\\
&&\qquad\qquad\qquad\qquad
\left.- \stackrel{\scriptscriptstyle _s}{\bmi T}{}_{\!\!\! x}(-k_y){\rf{\bmi 
G}}({\bmi x}_{R};{z, x};k_y)  
\stackrel{\scriptscriptstyle _r}{\bmi T}{}_{\!\!\! x}(k_y){{\bmi G}}({z, 
x};{\bmi x}_{S};k_y) \ {\rm d}z \right]
\end{eqnarray}
Using 
\[
\int_{{\mit \Omega}} \overbrace{{\bmi J}_{6} ({\bmi L}^{\dagger}{\rf{\bmi G}})^{T}(
{\bmi x};{\bmi x}_{R};-k_y){\bmi J}_{6}}^{{\bmi I}_{6}\delta({\bmi x}-{\bmi 
x}_R)}{{\bmi G}}({\bmi 
x};{\bmi x}_{S};k_y) {\rm d}{\bmi x}=\left\{
\begin{array}{cc}
{{\bmi G}}({\bmi x}_{R};{\bmi x}_{S};k_y)& {\bmi x}_{R}\in {\mit \Omega}\\
{\bf 0}&{\bmi x}_{R}\in {\mit \Omega}^{e}
\end{array}  \right.
\]
and
\begin{eqnarray*}
\lefteqn{
\int_{{\mit \Omega}} {\rf{\bmi G}}({\bmi x}_{R};{\bmi x};k_y)
\overbrace{{\bmi L}{{\bmi G}}({\bmi x};{\bmi x}_{S};k_y)}^{{\mit \Delta} {\bmi A}({\bmi
x},\partial_{x};k_y) {\bmi G}({\bmi x},{\bmi x}_{S};k_y)} {\rm d}{\bmi 
x}=}\\
&&\qquad\qquad\qquad\qquad\qquad
\int_{{\mit \Omega}} {\rf{\bmi G}}({\bmi x}_{R};{\bmi x};k_y)
{\mit \Delta} {\bmi A}({\bmi
x},\partial_{x};k_y) {\bmi G}({\bmi x},{\bmi x}_{S};k_y) {\rm d}{\bmi x}
\end{eqnarray*}
in (\ref{thm21}) and rearranging proves Thm.2.
\hfill \rule{11 pt}{11 pt}



%%%%%%%%%%%%%%%%%%%%%%%%%%%%%%%%%%%%%%%%%%%%%%%%%%%%%%%%%%%%%%%%%%%%%%%%%%%%%%%%%%%%%%%%%
%%%%%%%%%%%%%%%%%%%%%%%%%%%%%%%%%%%%%%%%%%%%%%%%%%%%%%%%%%%%%%%%%%%%%%%%%%%%%%%%%%%%%%%%%
%%%%%%%%%%%%%%%%%%%%%%%%%%%%%%%%%%%%%%%%%%%%%%%%%%%%%%%%%%%%%%%%%%%%%%%%%%%%%%%%%%%%%%%%%
\subsection{The Lippmann-Schwinger type integral equation} \label{secLS}

We can derive a number of useful identities using Thm.1 and Thm.2 by 
employing the {\em boundary integral\/} identity
\[
\left\langle{\rf{\bmi G}}({\bmi x};{\bmi x}_{1};-k_y),{{\bmi 
G}}({\bmi x};{\bmi x}_{2};k_y)\right\rangle_{\partial{\mit 
\Omega}^{-}}\equiv
-\left\langle{\rf{\bmi G}}({\bmi x};{\bmi x}_{1};-k_y),{{\bmi 
G}}({\bmi x};{\bmi x}_{2};k_y)\right\rangle_{\partial{\mit \Omega}^{+}}
\]
or equivalently
\begin{eqnarray} \label{bintid}
\lefteqn{
\oint_{\partial {\mit \Omega}^{-}} \left[ {\rf{\bmi G}}({\bmi x}_{R};{z, x};k_y) 
{{\bmi G}}({z, x};{\bmi x}_{S};k_y) \  {\rm d}x \right.}\nonumber\\
&&\qquad\qquad\qquad\qquad
\left.- \stackrel{\scriptscriptstyle _s}{\bmi T}{}_{\!\!\! x}(-k_y){\rf{\bmi 
G}}({\bmi x}_{R};{z, x};k_y)  
\stackrel{\scriptscriptstyle _r}{\bmi T}{}_{\!\!\! x}(k_y){{\bmi G}}({z, 
x};{\bmi x}_{S};k_y) \ {\rm d}z \right]\equiv\nonumber\\
&&
-\oint_{\partial {\mit \Omega}^{+}} \left[ {\rf{\bmi G}}({\bmi x}_{R};{z, x};k_y) 
{{\bmi G}}({z, x};{\bmi x}_{S};k_y) \  {\rm d}x \right.\nonumber\\
&&\qquad\qquad\qquad\qquad
\left.- \stackrel{\scriptscriptstyle _s}{\bmi T}{}_{\!\!\! x}(-k_y){\rf{\bmi 
G}}({\bmi x}_{R};{z, x};k_y)  
\stackrel{\scriptscriptstyle _r}{\bmi T}{}_{\!\!\! x}(k_y){{\bmi G}}({z, 
x};{\bmi x}_{S};k_y) \ {\rm d}z \right]
\end{eqnarray}
because the {\em sense of the contour integration has been reversed in 
the two integrals\/}.

First, let us assume that ${\bmi x}_{R}\in {\mit \Omega}$ in Thm.1 and Thm.2. From
Thm.1 it follows that
\begin{samepage}
\begin{eqnarray} \label{Huygens0}
\lefteqn{
{\rf{\bmi G}}({\bmi x}_{R};{\bmi x}_{S};k_y)=
\oint_{\partial {\mit \Omega}^{+}} \left[ {\rf{\bmi G}}({\bmi x}_{R};{z, x};k_y) 
{{\bmi G}}({z, x};{\bmi x}_{S};k_y) \  {\rm d}x \right.}\nonumber\\
&&
\qquad\qquad\qquad
- \left.\stackrel{\scriptscriptstyle _s}{\bmi T}{}_{\!\!\! x}(-k_y){\rf{\bmi 
G}}({\bmi x}_{R};{z, x};k_y)  
\stackrel{\scriptscriptstyle _r}{\bmi T}{}_{\!\!\! x}(k_y){{\bmi G}}({z, 
x};{\bmi x}_{S};k_y) \ {\rm d}z \right],\quad {\bmi x}_{R}\in {\mit 
\Omega}\nonumber\\
&&
\qquad\qquad\qquad
\oint_{\partial {\mit \Omega}^{+}} \left[ {\rf{\bmi G}}({\bmi x}_{R};{z, x};k_y) 
{{\bmi G}}({z, x};{\bmi x}_{S};k_y) \  {\rm d}x \right.\\
&&
\qquad\qquad\qquad
- \left.\stackrel{\scriptscriptstyle _s}{\bmi T}{}_{\!\!\! x}(-k_y){\rf{\bmi 
G}}({\bmi x}_{R};{z, x};k_y)  
\stackrel{\scriptscriptstyle _r}{\bmi T}{}_{\!\!\! x}(k_y){{\bmi G}}({z, 
x};{\bmi x}_{S};k_y) \ {\rm d}z \right],\quad {\bmi x}_{R}\in {\mit 
\Omega}\nonumber
\end{eqnarray}
\end{samepage}
where the last equation follows from (\ref{bintid}).
Similarly, setting 
${\bmi x}_{R}\in {\mit \Omega}$ in Thm.2., we get
\begin{eqnarray} \label{intrep1}
{{\bmi G}}({\bmi x}_{R};{\bmi x}_{S};k_y)
&=&\int_{{\mit \Omega}} {\rf{\bmi G}}({\bmi x}_{R};{\bmi x};k_y)
{\mit \Delta} {\bmi A}({\bmi
x},\partial_{x};k_y) {\bmi G}({\bmi x},{\bmi x}_{S};k_y) {\rm d}{\bmi x}\nonumber\\
&&
-\oint_{\partial {\mit \Omega}^{-}} \left[ {\rf{\bmi G}}({\bmi x}_{R};{z, x};k_y) 
{{\bmi G}}({z, x};{\bmi x}_{S};k_y) \  {\rm d}x \right.\\
&&
\qquad\qquad
- \left.\stackrel{\scriptscriptstyle _s}{\bmi T}{}_{\!\!\! x}(-k_y){\rf{\bmi 
G}}({\bmi x}_{R};{z, x};k_y)  
\stackrel{\scriptscriptstyle _r}{\bmi T}{}_{\!\!\! x}(k_y){{\bmi G}}({z, 
x};{\bmi x}_{S};k_y) \ {\rm d}z \right]\nonumber
\end{eqnarray}

Now, using (\ref{Huygens0}) in (\ref{intrep1}), we get the {\em interior 
Lippmann-Schwinger type integral equation\/}
\begin{samepage}
\begin{eqnarray} \label{LSinterior}
{{\bmi G}}({\bmi x}_{R};{\bmi x}_{S};k_y)&=&
{\rf{\bmi G}}({\bmi x}_{R};{\bmi x}_{S};k_y)\nonumber\\
&+&\int_{{\mit \Omega}} {\rf{\bmi G}}({\bmi x}_{R};{\bmi x};k_y)
{\mit \Delta} {\bmi A}({\bmi
x},\partial_{x};k_y) {\bmi G}({\bmi x},{\bmi x}_{S};k_y) {\rm d}{\bmi x},
\quad {\bmi x}_{R}\in {\mit \Omega}
\end{eqnarray}
\end{samepage}

Next, we set ${\bmi x}_{R}\in {\mit \Omega}^{e}$ in Thm1. to get
\begin{eqnarray} \label{extrep1}
\lefteqn{
{{\bmi G}}({\bmi x}_{R};{\bmi x}_{S};k_y)
-{\rf{\bmi G}}({\bmi x}_{R};{\bmi x}_{S};k_y)=
}\\
&&
\qquad\qquad
-\oint_{\partial {\mit \Omega}^{+}} \left[ {\rf{\bmi G}}({\bmi x}_{R};{z, x};k_y) 
{{\bmi G}}({z, x};{\bmi x}_{S};k_y) \  {\rm d}x \right.\nonumber\\
&&
\qquad\qquad\qquad\qquad
- \left.\stackrel{\scriptscriptstyle _s}{\bmi T}{}_{\!\!\! x}(-k_y){\rf{\bmi 
G}}({\bmi x}_{R};{z, x};k_y)  
\stackrel{\scriptscriptstyle _r}{\bmi T}{}_{\!\!\! x}(k_y){{\bmi G}}({z, 
x};{\bmi x}_{S};k_y) \ {\rm d}z \right]\nonumber\\
&&
\qquad\qquad
=\oint_{\partial {\mit \Omega}^{-}} \left[ {\rf{\bmi G}}({\bmi x}_{R};{z, x};k_y) 
{{\bmi G}}({z, x};{\bmi x}_{S};k_y) \  {\rm d}x \right.\nonumber\\
&&
\qquad\qquad
- \left.\stackrel{\scriptscriptstyle _s}{\bmi T}{}_{\!\!\! x}(-k_y){\rf{\bmi 
G}}({\bmi x}_{R};{z, x};k_y)  
\stackrel{\scriptscriptstyle _r}{\bmi T}{}_{\!\!\! x}(k_y){{\bmi G}}({z, 
x};{\bmi x}_{S};k_y) \ {\rm d}z \right], \quad {\bmi x}_{R}\in {\mit \Omega}^{e}\nonumber
\end{eqnarray}
where the last equation in (\ref{extrep1}) follows from (\ref{bintid}). 
Similarly,  we set ${\bmi x}_{R}\in {\mit \Omega}^{e}$ in Thm2. to get
\begin{eqnarray} \label{extrep2}
\lefteqn{\int_{{\mit \Omega}} {\rf{\bmi G}}({\bmi x}_{R};{\bmi x};k_y)
{\mit \Delta} {\bmi A}({\bmi
x},\partial_{x};k_y) {\bmi G}({\bmi x},{\bmi x}_{S};k_y) {\rm d}{\bmi x}=}\\
&&
\oint_{\partial {\mit \Omega}^{-}} \left[ {\rf{\bmi G}}({\bmi x}_{R};{z, x};k_y) 
{{\bmi G}}({z, x};{\bmi x}_{S};k_y) \  {\rm d}x \right.\nonumber\\
&&
\qquad\qquad
- \left.\stackrel{\scriptscriptstyle _s}{\bmi T}{}_{\!\!\! x}(-k_y){\rf{\bmi 
G}}({\bmi x}_{R};{z, x};k_y)  
\stackrel{\scriptscriptstyle _r}{\bmi T}{}_{\!\!\! x}(k_y){{\bmi G}}({z, 
x};{\bmi x}_{S};k_y) \ {\rm d}z \right], \quad {\bmi x}_{R}\in {\mit \Omega}^{e}\nonumber
\end{eqnarray}
Using the domain integral on the left hand side of (\ref{extrep2}) in 
place of the boundary integral on the right hand side of (\ref{extrep1}) 
we get the {\em exterior Lippmann-Schwinger type integral equation\/}
\begin{eqnarray} \label{LSexterior}
{{\bmi G}}({\bmi x}_{R};{\bmi x}_{S};k_y)&=&
{\rf{\bmi G}}({\bmi x}_{R};{\bmi x}_{S};k_y)\nonumber\\
&+&\int_{{\mit \Omega}} {\rf{\bmi G}}({\bmi x}_{R};{\bmi x};k_y)
{\mit \Delta} {\bmi A}({\bmi
x},\partial_{x};k_y) {\bmi G}({\bmi x},{\bmi x}_{S};k_y) {\rm d}{\bmi x},
\quad {\bmi x}_{R}\in {\mit \Omega}^{e}
\end{eqnarray}


{\noindent{\bf Corollary 1:}{\em The Domain Integral 
Representation\/}}\newline
Given that
 \begin{eqnarray*}
\lefteqn{{\bmi L}({\bmi x},\partial_{z},\partial_{x};k_y) {\bmi G}({\bmi
x},{\bmi x}_{S};k_y) =}\\
&& \qquad {\bmi I}_{6}\delta({\bmi x}-{\bmi x}_{S}) + 
{\chi}_{{{\mit \Omega}}} {\mit \Delta} {\bmi A}({\bmi
x},\partial_{x};k_y) {\bmi G}({\bmi x},{\bmi x}_{S};k_y),\qquad {\bmi 
x}_{S}\in{\mit \Omega}^{e}
\end{eqnarray*}
and
\begin{displaymath}
{\bmi L}^{\dagger}({\bmi x},\partial_{z},\partial_{x};k_y) \rf{\bmi G}({\bmi
x},{\bmi x}_{R};-k_y) = -{\bmi I}_{6}\delta({\bmi x}-{\bmi x}_{R}),\ \ \ \ {\bmi 
x}_{R}\in{R}^{2}
\end{displaymath}
then
\begin{eqnarray*}
{{\bmi G}}({\bmi x}_{R};{\bmi x}_{S};k_y)&=&
{\rf{\bmi G}}({\bmi x}_{R};{\bmi x}_{S};k_y)\\
&&+\int_{{\mit \Omega}} {\rf{\bmi G}}({\bmi x}_{R};{\bmi x};k_y)
{\mit \Delta} {\bmi A}({\bmi
x},\partial_{x};k_y) {\bmi G}({\bmi x},{\bmi x}_{S};k_y) {\rm d}{\bmi x},
\quad {\bmi x}_{R}\in R^{2}
\end{eqnarray*}
{\em Proof\/} Follows directly from (\ref{LSinterior}) and 
(\ref{LSexterior}).
\hfill \rule{11 pt}{11 pt}



 In words, 
{\bf Corollary 1} is a statement that:
 the $6 \times 6$ Green's tensor  ${{\bmi G}}({\bmi x}_{R};{\bmi 
x}_{S};k_y)$ for a heterogeneous medium is the sum of the incident 
Green's tensor ${\rf{\bmi G}}({\bmi x}_{R};{\bmi x}_{S};k_y)$ propagating 
through the reference medium
and the radiated (or scattered) field $\int_{{\mit \Omega}} {\rf{\bmi G}}({\bmi x}_{R};{\bmi x};k_y)
{\mit \Delta} {\bmi A}({\bmi
x},\partial_{x};k_y) {\bmi G}({\bmi x},{\bmi x}_{S};k_y) {\rm d}{\bmi x}$ 
from the {\em secondary source distribution\/} ${\mit \Delta} {\bmi A}({\bmi
x},\partial_{x};k_y) {\bmi G}({\bmi x},{\bmi x}_{S};k_y)$,  supported on 
${\mit \Omega}$. In general, the larger the {\em contrast\/} ${\mit \Delta} {\bmi A}({\bmi
x},\partial_{x};k_y)$ between the actual medium, and the reference medium, 
the more difficult it will be to find the solution 
${{\bmi G}}({\bmi x}_{R};{\bmi x}_{S};k_y)$. The rest of the thesis will be concerned with solving this 
Lipmann-Schwinger type operator equation.













